\documentclass[12pt, a4paper]{article}

\usepackage[utf8]{inputenc}
\usepackage[T2A]{fontenc}
\usepackage[russian]{babel}
\usepackage{amsmath, amssymb, amsthm}
\usepackage{geometry}
\usepackage{hyperref}
\geometry{top=2cm, bottom=2cm, left=2.5cm, right=1.5cm}

\newtheorem{theorem}{Теорема}
\newtheorem{lemma}{Лемма}
\newtheorem{definition}{Определение}
\newtheorem{example}{Пример}

\title{Билет №9 \\ \small Алгоритмы на графах. Обходы графов. Кратчайшие пути, семейство алгоритмов A*. Остовные деревья. Задача о максимальном потоке, о паросочетании, о потоке минимальной стоимости. (ИТМО) \\ Примеры эффективных (полиномиальных) алгоритмов: быстрые алгоритмы поиска и сортировки; полиномиальные алгоритмы для задач на графах и сетях (поиск в глубину и ширину, о минимальном остове, о кратчайшем пути, о назначениях). (СПбГУ)
}
\author{Сериков Владислав Александрович}
\date{20 апреля 2026}

\begin{document}

\maketitle
\tableofcontents
\newpage

\section{Примеры эффективных (полиномиальных) алгоритмов}

\begin{definition}
\textbf{Полиномиальный алгоритм} --- это алгоритм, время работы которого ограничено сверху многочленом от размера входных данных: $O(n^k)$, где $k$ --- константа, $n$ --- размер входа.
\end{definition}

К классу полиномиальных алгоритмов относятся быстрые алгоритмы поиска и сортировки, а также многие алгоритмы на графах (поиск пути, остовные деревья и др.). Задачи, решаемые такими алгоритмами, образуют класс $P$.

\subsection{Быстрые алгоритмы поиска и сортировки}
\textbf{Бинарный поиск (Binary Search).} Ищет элемент в отсортированном массиве за $O(\log n)$.
\begin{itemize}
    \item \textit{Шаги:} 
    \begin{itemize}
        \item Выбрать средний элемент. 
        \item Если он равен искомому --- конец. 
        \item Если он больше, рекурсивно искать в левой половине, иначе в правой.
    \end{itemize}
    \item \textit{Пример:} Ищем 5 в массиве $[1, 3, 5, 7]$. Средний элемент 3 (индекс 1). $3 < 5$, ищем в $[5, 7]$. Средний 5. Нашли!
\end{itemize}

\textbf{Сортировка слиянием (Merge Sort).} Сортирует массив за $O(n \log n)$.
\begin{itemize}
    \item \textit{Шаги:} 
    \begin{itemize}
        \item Разделить массив пополам. 
        \item Рекурсивно отсортировать половины. 
        \item Слить два отсортированных массива за $O(n)$.
    \end{itemize}
    \item \textit{Пример:} $[4, 1, 3, 2] \to [4, 1]$ и $[3, 2] \to[1, 4]$ и $[2, 3]$. Слияние: берем меньшие элементы $\to[1, 2, 3, 4]$.
\end{itemize}

\textbf{Быстрая сортировка (Quick Sort).} Эффективная сортировка in-place на основе парадигмы «разделяй и властвуй». Среднее время $O(n \log n)$, в худшем случае $O(n^2)$.
\begin{itemize}
    \item \textit{Шаги:} 
    \begin{itemize}
        \item Выбрать опорный элемент (pivot). 
        \item Разбить массив так, чтобы слева от опорного были элементы меньше него, а справа --- больше или равные (Partition). 
        \item Рекурсивно отсортировать левую и правую части.
    \end{itemize}
    \item \textit{Пример:} Массив $[3, 1, 4, 2, 5]$. Выберем опорным $3$. Разделение: слева элементы $<3$ $\to[1, 2]$, справа $\ge 3 \to [4, 5]$. Сам опорный встает между ними: $[1, 2] \cup \{3\} \cup [4, 5] \to[1, 2, 3, 4, 5]$.
\end{itemize}

\textbf{Сортировка пузырьком (Bubble Sort).} Простая сортировка обменом. Время $O(n^2)$ в среднем и худшем случае, $O(n)$ в лучшем (если массив уже отсортирован).
\begin{itemize}
    \item \textit{Шаги:}
    \begin{itemize}
        \item Пройти по массиву от начала до конца. 
        \item Сравнивать соседние элементы: если левый больше правого, поменять их местами. 
        \item Повторять проходы, пока на очередном проходе не будет ни одного обмена. На каждом шаге наибольший элемент «всплывает» в конец.
    \end{itemize}
    \item \textit{Пример:} Массив $[3, 2, 1]$. Шаг 1: $3 > 2 \implies [2, 3, 1]$. Шаг 2: $3 > 1 \implies [2, 1, 3]$. Максимум (3) на месте. Шаг 3: $2 > 1 \implies [1, 2, 3]$.
\end{itemize}

\textbf{Сортировка вставками (Insertion Sort).} Эффективна для небольших или частично отсортированных массивов (online алгоритм). Время $O(n^2)$, лучшая оценка $O(n)$.
\begin{itemize}
    \item \textit{Шаги:}
    \begin{itemize}
        \item Считать первый элемент отсортированным подмассивом. 
        \item Взять следующий элемент и сдвинуть все бОльшие элементы отсортированной части вправо. 
        \item Вставить взятый элемент на освободившееся место. 
        \item Повторить для всех элементов.
    \end{itemize}
    \item \textit{Пример:} Массив $[4, 2, 3]$.
    \begin{itemize}
        \item \underline{4} отсортировано.
        \item Берем $2$: сдвигаем $4$, вставляем $2 \to [\underline{2, 4}, 3]$. 
        \item Берем $3$: сдвигаем $4$, вставляем $3 \to [\underline{2, 3, 4}]$.
    \end{itemize}
\end{itemize}

\textbf{Сортировка подсчётом (Counting Sort).} Не основана на сравнениях (позволяет обойти нижнюю оценку $\Omega(n \log n)$ для сортировок сравнениями). Работает с целыми числами. Время и память: $O(n + k)$, где $k$ --- диапазон значений.
\begin{itemize}
    \item \textit{Шаги:}
    \begin{itemize}
        \item Найти максимальное значение в массиве ($k$). 
        \item Создать массив частот размера $k+1$ и подсчитать количество вхождений каждого элемента. 
        \item Последовательно пройти по массиву частот и перезаписать исходный массив в соответствии с посчитанными количествами (для сохранения стабильности вычисляют префиксные суммы).
    \end{itemize}
    \item \textit{Пример:} Массив $[2, 0, 2, 1]$. Диапазон $k=2$. Массив частот для $0, 1, 2$: $C[0]=1, C[1]=1, C[2]=2$. Перезапись: один `0`, один `1`, два `2` $\to[0, 1, 2, 2]$.
\end{itemize}

\section{Алгоритмы на графах: Обходы и кратчайшие пути}

\subsection{Обходы графов (BFS и DFS)}
\textbf{Поиск в ширину (BFS).} Находит кратчайшие пути (по числу ребер) от стартовой вершины в невзвешенном графе. Время работы: $O(V + E)$.
\begin{itemize}
    \item \textit{Шаги:} 
    \begin{itemize}
        \item Поместить стартовую вершину в очередь и отметить как посещенную. 
        \item Пока очередь не пуста, извлечь вершину $u$. 
        \item Всех непосещенных соседей $u$ отметить и добавить в очередь.
    \end{itemize}
    \item \textit{Пример:} Граф: $1-2, 1-3, 2-4$. Старт: 1. Очередь: $[1]$. Извлекаем 1, добавляем 2, 3. Очередь: $[2, 3]$. Извлекаем 2, добавляем 4. Порядок обхода: 1, 2, 3, 4.
\end{itemize}

\textbf{Поиск в глубину (DFS).} Используется для топологической сортировки, поиска компонент сильной связности. Время работы: $O(V + E)$.
\begin{itemize}
    \item \textit{Шаги:}
    \begin{itemize}
        \item Поместить стартовую вершину в стек (или использовать рекурсию). 
        \item Пометить как посещенную. 
        \item Рекурсивно вызвать DFS для первого непосещенного соседа. Если таких нет --- вернуться назад (backtrack).
    \end{itemize}
\end{itemize}

\subsection{Кратчайшие пути и семейство алгоритмов A*}
Для взвешенных графов (без отрицательных весов) используется \textbf{Алгоритм Дейкстры} ($O(E \log V)$ с мин-кучей). Он последовательно релаксирует ребра из вершины с минимальным текущим расстоянием.

\textbf{Алгоритм A* (A-star)} --- это информированный алгоритм поиска, обобщающий Дейкстру за счет использования эвристики.
\begin{itemize}
    \item \textit{Функция оценки:} $f(v) = g(v) + h(v)$, где $g(v)$ --- точная стоимость пути от старта до $v$, $h(v)$ --- эвристическая оценка стоимости от $v$ до цели.
    \item \textit{Допустимость (Admissibility):} Если $h(v)$ никогда не переоценивает реальное расстояние до цели (т.е. $h(v) \le h^*(v)$), то A* гарантированно находит кратчайший путь.
    \item \textit{Шаги:} 
    \begin{itemize}
        \item Поместить старт в очередь (Open Set) с приоритетом $f(s) = h(s)$. 
        \item Пока очередь не пуста, извлечь $u$ с мин. $f(u)$. 
        \item Если $u$ --- цель, завершить. 
        \item Иначе обновить $g(v)$ для соседей и добавить их в очередь.
    \end{itemize}
\end{itemize}

\section{Остовные деревья}

\begin{definition}
\textbf{Остовное дерево (Spanning Tree)} связного неориентированного графа $G$ --- это подграф, являющийся деревом и содержащий все вершины графа $G$.
\end{definition}

\textbf{Количество остовных деревьев.}
\begin{itemize}
    \item \textbf{Формула Кэли (Cayley's formula):} Количество различных остовных деревьев в полном графе $K_n$ равно $n^{n-2}$.
    \item \textbf{Теорема Татта — Нэша-Уильямса:} Граф содержит $k$ попарно непересекающихся (по ребрам) остовных деревьев тогда и только тогда, когда для любого разбиения множества вершин на $r$ компонент количество ребер между этими компонентами не меньше $k(r-1)$.
\end{itemize}

\begin{definition}
\textbf{Минимальное остовное дерево (MST)} --- остовное дерево, сумма весов ребер которого минимальна.
\end{definition}

\textbf{Алгоритм Прима} (Полиномиальный алгоритм, $O(E \log V)$):
\begin{itemize}
    \item \textit{Шаги:}
    \begin{itemize}
        \item Начать с любой вершины, добавить её в растущее дерево $T$.
        \item На каждом шаге выбирать ребро минимального веса, соединяющее вершину из $T$ с вершиной не из $T$, и добавлять его в $T$.
        \item Повторять до охвата всех вершин.
    \end{itemize}
    \item \textit{Доказательство корректности (свойство разреза):} Любой разрез графа, разделяющий $T$ и остальные вершины, должен быть пересечен хотя бы одним ребром. Ребро минимального веса в этом разрезе обязано принадлежать минимальному остову.
\end{itemize}

\section{Сети, Потоки и Разрезы}

\begin{definition}
\textbf{Граф-сеть} --- ориентированный граф $G=(V, E)$ с выделенным истоком $s$ и стоком $t$, где каждому ребру $e$ приписана пропускная способность $c(e) \ge 0$.
\end{definition}

\begin{definition}
\textbf{Поток} --- это функция $f: E \to \mathbb{R}$, удовлетворяющая условиям:
\begin{enumerate}
    \item Ограничение пропускной способности: $0 \le f(e) \le c(e)$ (здесь 0 --- \textbf{нижняя оценка}, $c(e)$ --- \textbf{верхняя оценка} потока).
    \item Сохранение потока: для любой вершины $v \neq s, t$, $\sum_{u} f(u, v) = \sum_{w} f(v, w)$.
\end{enumerate}
\end{definition}

\begin{definition}
\textbf{Остаточная сеть $G_f$} для потока $f$ --- граф с \textbf{остаточными пропускными способностями} $c_f(u, v) = c(u, v) - f(u, v) + f(v, u)$. 
Путь из $s$ в $t$ в остаточной сети называется \textbf{дополняющим (остаточным) путем}.
\end{definition}

\begin{definition}
\textbf{Разрез $(S, T)$} --- разбиение $V$ на $S$ и $T$, где $s \in S, t \in T$. Пропускная способность разреза: $c(S, T) = \sum_{u \in S, v \in T} c(u, v)$.
\end{definition}

\begin{theorem}
\textbf{Теорема Форда-Фалкерсона (Min-cut Max-flow):} Величина максимального потока равна пропускной способности минимального разреза.
\end{theorem}
\begin{proof}
(Идея) Очевидно, что любой поток $|f|$ не превышает емкость любого разреза $c(S, T)$ (Слабая двойственность). Алгоритм Форда-Фалкерсона ищет дополняющие пути. Когда таких путей больше нет, множество достижимых из $s$ вершин в $G_f$ образует множество $S$, а остальные --- $T$. Для всех ребер из $S$ в $T$ $f(u, v) = c(u, v)$, а из $T$ в $S$ поток равен 0. Следовательно, $|f| = c(S, T)$, что доказывает как максимальность потока, так и минимальность данного разреза.
\end{proof}

\subsection{Глобальный минимальный разрез: Алгоритм Каргера}
Находит минимальный разрез в неориентированном графе без выделенных истока и стока (вероятностный полиномиальный алгоритм).
\begin{itemize}
    \item \textit{Шаги:}
    \begin{itemize}
        \item Выбрать случайное ребро.
        \item Стянуть его (объединить концы в одну вершину, удалить петли).
        \item Повторять, пока не останется 2 вершины. Оставшиеся ребра между ними образуют разрез.
    \end{itemize}
    \item \textit{Корректность:} Вероятность того, что минимальный разрез не будет уничтожен за одну итерацию алгоритма, достаточно велика, чтобы при полиномиальном числе перезапусков найти точный ответ с вероятностью, близкой к 1.
\end{itemize}

\subsection{Поток минимальной стоимости}
Каждому ребру, помимо пропускной способности $c(e)$, задана стоимость $a(e)$. Задача: найти поток заданной величины (или максимальный), суммарная стоимость которого $\sum f(e)a(e)$ минимальна.
\begin{itemize}
    \item \textit{Алгоритм (Successive shortest path):} Ищем дополняющий путь минимальной стоимости в $G_f$ (с помощью алгоритма Форда-Беллмана или Дейкстры с потенциалами) и пускаем по нему максимально возможный поток.
\end{itemize}

\section{Двудольные графы и Паросочетания}

\begin{definition}
\textbf{Двудольный граф} --- граф, множество вершин которого можно разбить на две непересекающиеся доли так, что каждое ребро соединяет вершины из разных долей.
\end{definition}

\textbf{Задача о трёх домиках и трёх колодцах:} Можно ли проложить дорожки от 3 домиков к 3 колодцам (газ, вода, электричество) так, чтобы они не пересекались?
\begin{proof}
Эта задача эквивалентна укладке графа $K_{3,3}$ на плоскости. По формуле Эйлера для планарного графа $V - E + F = 2$. У нас $V=6, E=9 \implies F=5$. Так как граф двудольный, в нем нет циклов длины 3, минимальный цикл $\ge 4$. Значит, $4F \le 2E \implies 20 \le 18$. Противоречие. Граф не планарный.
\end{proof}

\begin{definition}
\textbf{Паросочетание} --- множество попарно несмежных ребер.
\begin{itemize}
    \item \textbf{Максимальное} --- паросочетание с наибольшим числом ребер.
    \item \textbf{Минимальное} --- обычно понимают как минимальное по весу совершенное паросочетание (задача о назначениях) или минимальное максимальное паросочетание (наименьшее по размеру, к которому нельзя добавить ребро).
    \item \textbf{Совершенное} --- паросочетание, покрывающее все вершины графа.
\end{itemize}
\textbf{Вершинное покрытие} --- множество вершин, инцидентное каждому ребру графа. (По теореме Кёнига в двудольных графах размер макс. паросочетания равен размеру мин. вершинного покрытия).
\end{definition}

\begin{theorem}
\textbf{Теорема Холла:} В двудольном графе с долями $L$ и $R$ существует паросочетание, покрывающее долю $L$, тогда и только тогда, когда для любого подмножества вершин $S \subseteq L$ размер их окрестности $|N(S)| \ge |S|$.
\end{theorem}
\begin{proof}
(Идея) Необходимость очевидна. Достаточность легко доказывается через теорему Форда-Фалкерсона. Строим сеть: исток $s$ соединен с $L$, $R$ соединен со стоком $t$ (емкости 1), между $L$ и $R$ емкости $\infty$. Если максимальный поток меньше $|L|$, то минимальный разрез имеет емкость $< |L|$. Этот разрез выделяет подмножество $S \subseteq L$, для которого $|N(S)| < |S|$, что противоречит условию.
\end{proof}

\subsection{Венгерский алгоритм}
Решает задачу о назначениях (поиск совершенного паросочетания минимального веса в полном двудольном графе) за $O(V^3)$. Это строгий полиномиальный алгоритм.
\begin{itemize}
    \item \textit{Шаги:} 
    \begin{itemize}
        \item В матрице стоимостей вычесть из каждой строки ее минимум.
        \item Вычесть из каждого столбца его минимум.
        \item Попытаться найти совершенное паросочетание из нулевых элементов (найти независимые нули). Если оно найдено — это оптимальное назначение.
        \item Если нет, провести минимальное число линий через строки и столбцы, чтобы покрыть все нули.
        \item Найти минимальный непокрытый элемент. Вычесть его из всех непокрытых элементов и добавить к элементам на пересечении линий. Вернуться к шагу 3.
    \end{itemize}
\end{itemize}

\end{document}